\chapter{How to Read This Book}

This book teaches every technology inside the TeacherSupporter project ---
Spring Boot and Spring Cloud, Postgres and MongoDB, Kafka, Docker,
Kubernetes, and the Jenkins pipeline that ties them together --- using the
project's own files as the worked examples. Nothing is invented for the
page: every configuration you will read is deployed, or about to be
deployed, on a real two-core home server called \code{ts-server}.

\section*{The rules the book follows}

\begin{description}
  \item[Concepts before configs.] Each chapter first teaches the idea in
  vendor-neutral terms, then reads the project's real file line by line,
  then shows one or two examples from \emph{outside} the project so you can
  tell what generalizes from what is local habit.
  \item[Decisions are first-class.] Wherever the project chose between
  alternatives (Jib over \code{docker build}, Kustomize over Helm, a second
  Postgres server over a second database) the reasoning is in a
  \textsc{decision} box. These are the ``why did you\ldots'' questions
  interviews are made of.
  \item[Failure is shown symptom-first.] \textsc{pitfall} boxes start with
  what you would actually see (\code{CrashLoopBackOff}, a TLS error, a
  silent 404) and work back to the cause, because that is the direction
  real debugging runs in.
\end{description}

\section*{The box legend}

\begin{decision}[example]
Why the project chose one design over another, and what would have to
change for the other choice to win.
\end{decision}

\begin{pitfall}[example]
A symptom you will eventually meet, followed by its cause and fix.
\end{pitfall}

\begin{handson}[example]
A command to run against the real stack, with the output you should
expect. The book is meant to be read next to a terminal.
\end{handson}

\begin{hardware}[example]
A choice that is only correct on this machine (2 cores, 19.4\,GB RAM,
SATA SSD, Wi-Fi). Do not copy these into a data center.
\end{hardware}

\begin{quiz}
Every chapter ends with a few questions in this box. Answers are in
Appendix~A. Write your answer down before checking --- recognition is not
recall.
\end{quiz}

\section*{Notation}

Inline code and file names look like \code{application.yml}. Annotated
listings carry circled markers --- \co{1}, \co{2} --- explained directly
beneath the listing. Shell commands are shown with a \code{\$} prompt for
the server and \code{PS>} for Windows PowerShell; do not type the prompt.

\section*{What is assumed, and what is not}

Assumed: you can read Java, and you have used Maven and git. Not assumed:
any prior Kubernetes, Jenkins, Kafka, Docker internals, or frontend
knowledge. Where the project's own documentation already covers one-time
machine setup (\code{docs/RUNBOOK-SERVER-SETUP.md}), this book summarizes
and moves on.
